\subsection*{Chapter 15: Sharded Data Parallelism from Scratch}

\paragraph{Exercise 15.1}
Algorithm bandwidth is $n/t$. A bandwidth-bound ring all-reduce takes $t = 2\frac{p-1}{p}\frac{n}{W}$,
so $\frac{n}{t}\cdot\frac{2(p-1)}{p} = W$: the bus bandwidth equals the link bandwidth for every $p$,
which is what makes it comparable across world sizes. An all-gather of a tensor whose full size is $n$ takes
$t = \frac{p-1}{p}\frac{n}{W}$ (each rank receives every chunk but its own once), so the matching
correction factor is $\frac{p-1}{p}$, and again bus bandwidth $=W$.

\paragraph{Exercise 15.2}
Gradients appear at rate $G/T_b$, so a bucket of $b$ bytes is filled every $bT_b/G$ seconds and costs
$A(b) = 2(p-1)\alpha + 2\frac{p-1}{p}b\beta$ to all-reduce. If $A(b) \le bT_b/G$, communication keeps
pace with backward and only the last bucket is exposed, costing $A(b)$, which grows with $b$. If
$A(b) > bT_b/G$, a queue builds and the exposed time is roughly $\frac{G}{b}A(b) - T_b$, which
\emph{falls} with $b$. The minimum is therefore at the smallest bucket that keeps pace:
\[
  b^\star = \frac{2(p-1)\alpha}{T_b/G - 2\frac{p-1}{p}\beta}.
\]
With $G=1.2$\,GB, $T_b=\SI{50}{\milli\second}$, $p=8$, $\alpha=\SI{10}{\micro\second}$ and 450\,GB/s:
$T_b/G = 4.17\times10^{-11}$\,s/B and $2\frac78\beta = 3.89\times10^{-12}$\,s/B, so
$b^\star \approx 1.4\times10^{-4}/3.78\times10^{-11}\approx3.7$\,MB, and the exposed tail is
$A(b^\star)\approx\SI{0.15}{\milli\second}$. In practice one uses a somewhat larger bucket to absorb
uneven gradient arrival, and a small first bucket.

\paragraph{Exercise 15.3}
After the reduce-scatter each rank holds averaged gradients for its shard only. Compute the local
squared norm $s_r = \lVert g^{(r)}\rVert_2^2$ (padding elements are zero, so they do not contribute),
all-reduce the scalar to get $\sum_r s_r$, take the square root, and scale the shard gradient by
$\min(1, c/(\lVert g\rVert + \epsilon))$. That is one scalar all-reduce, 4--8 bytes and pure
latency. The test spawns 3 ranks, trains one step on a model whose gradients exceed the clip threshold
(scale the loss up), and compares the updated parameters with a single process that calls
\code{torch.nn.utils.clip\_grad\_norm\_} before \code{AdamW.step()}, to $10^{-6}$.

\paragraph{Exercise 15.4}
Save, per parameter group, the parameter shapes, the unpadded flat length $N$, the old world size $p$
(hence shard size $s=\lceil N/p\rceil$), and each rank's shard of the master weights and of every
optimizer-state tensor with the same layout (Adam's $m$, $v$), plus replicated scalars such as the step
count. New rank $r'$ with shard size $s'=\lceil N/p'\rceil$ owns flat indices
$[r's', \min((r'+1)s', N))$ and reads the old shards $\lfloor r's'/s\rfloor$ through
$\lfloor((r'+1)s'-1)/s\rfloor$, slicing the overlapping ranges and zero-padding the tail. In the
parameter-ownership layout of \cref{ch:distributed}, state is stored per whole tensor on its owner. It can
be reloaded on $p'$ ranks by recomputing ownership with \code{partition\_parameters} and reading each
tensor's state from whichever old file holds it, which is possible but needs a name-to-owner map, and
balance is again limited by the largest tensor.

\paragraph{Exercise 15.5}
\code{self.flat.copy\_(buf)} is an in-place op on \code{flat}. Views and \code{detach()} share a version
counter with their base, so the counter of every view autograd saved during forward (each weight used
by a linear layer) increments. In backward, autograd compares each saved tensor's version with the
version recorded at save time, finds a mismatch, and raises \emph{``one of the variables needed for
gradient computation has been modified by an inplace operation: \dots\ is at version $k+1$; expected
version $k$''}. A storage-level copy changes the bytes without touching version counters. That is
safe here because the bytes written are exactly the bytes that were there when the tensors were saved,
the same gathered shards with no optimizer step in between. It would be unsafe if the weights could
change between forward and backward.

\paragraph{Exercise 15.6}
With $k$ blocks per unit there are $L/k$ units. Transient memory is about two gathered units plus one
unit's full gradient, $3kP_b\times2$ bytes in bf16. Collective latency per step is
$3\cdot\frac{L}{k}\cdot(p-1)\alpha$ (two all-gathers and one reduce-scatter per unit). Step time
falls with $k$ until the transient memory hits the budget $M$, so $k^\star = \lfloor M/(6P_b)\rfloor$
(capped by $L$). Larger units also mean coarser prefetch, but that matters only when $T/n$ is
near $C/W$. Qwen3-1.7B has $P_b \approx 5.03\times10^7$ ($\dmodel=2048$, 16 query and 8 KV heads of
128, MLP width 6144), so a unit costs $\approx0.30k$\,GB of transient memory. On a 16\,GB GPU with $p=8$
the persistent state is $16\Psi/8\approx3.4$\,GB and the root unit (tied 151{,}936$\times$2048
embedding) needs about 1.25\,GB gathered with its gradient. Leaving 6--8\,GB for activations, about
4\,GB remain for units, so $k\approx13$. Latency then falls from $28\cdot3\cdot7\cdot\SI{10}{\micro\second}
\approx\SI{5.9}{\milli\second}$ at $k=1$ to about $\SI{0.5}{\milli\second}$. One block per unit
is usually fine whenever per-step compute is tens of milliseconds.

\paragraph{Exercise 15.7}
$P = 8D^2$ (two matrices of $4D^2$). Per device, forward is $16TD^2/n$ FLOPs and backward $32TD^2/n$.
\begin{itemize}
  \item \textbf{DP}: forward has no communication. Backward all-reduces $16D^2$ bytes:
    $t_\text{comm} = 2\frac{n-1}{n}\frac{16D^2}{W}$ against $t_\text{comp}=\frac{32TD^2}{nC}$.
    It is compute-bound iff $n-1 \le TW/C$, that is $T/n \gtrsim C/W$.
  \item \textbf{FSDP}: forward all-gathers $\frac{n-1}{n}\frac{16D^2}{W}$ against $\frac{16TD^2}{nC}$, giving
    $n-1\le TW/C$. Backward all-gathers again and reduce-scatters, $2\frac{n-1}{n}\frac{16D^2}{W}$ against
    $\frac{32TD^2}{nC}$, with the same condition. FSDP scales as far as DP.
  \item \textbf{TP} (column-parallel $W_1$, row-parallel $W_2$): forward all-reduces the $(T,D)$ output,
    $2\frac{n-1}{n}\frac{2TD}{W}$ against $\frac{16TD^2}{nC}$, giving $n-1 \le 4DW/C$. Backward
    all-reduces the $(T,D)$ input gradient, the same bytes against twice the compute, giving
    $n-1\le8DW/C$. Forward binds: $n_\text{max} = 1 + 4DW/C$, independent of $T$. For $D=4096$ on
    H100/NVLink ($C/W\approx2200$) this is $\approx8.4$.
\end{itemize}

\paragraph{Exercise 15.8}
Each device holds $W_1, W_2$ sharded over both axes, $8D^2/(n_fn_t)$ elements, and $T/n_f$ tokens.
Forward: an all-gather over the \code{fsdp} axis of the device's TP slice,
$\frac{n_f-1}{n_f}\frac{16D^2/n_t}{W_f}$, against compute $\frac{16TD^2}{n_fn_tC}$, giving
$n_f-1\le TW_f/C$. Then TP's all-reduce over the \code{tp} axis of a $(T/n_f, D)$ activation,
$2\frac{n_t-1}{n_t}\frac{2TD/n_f}{W_t}$, against the same compute, gives $n_t-1\le 4DW_t/C$. The backward
pass is analogous (all-gather and reduce-scatter on \code{fsdp}, one all-reduce on \code{tp}). The two
conditions involve different variables, so they decouple. The coupling terms are: (i) the global
batch $T$ is shared, and $T/n_f$ must also stay large enough for efficient matmuls and for TP's
per-token compute; (ii) if both axes share physical links, $W_f$ and $W_t$ are not independent;
(iii) latency terms ($\alpha$) grow with the number of collectives, which both degrees increase;
(iv) activation memory per device scales with $T/n_f$.

\paragraph{Exercise 15.9}
\[
T_\text{hier} = 2\Big[(g-1)\alpha_1 + \tfrac{g-1}{g}n\beta_1\Big] + 2(N-1)\alpha_2 + 2\tfrac{N-1}{N}\tfrac{n}{g}\beta_2,
\qquad
T_\text{flat} = 2(gN-1)\alpha_2 + 2\tfrac{gN-1}{gN}n\beta_2 .
\]
With $\alpha_1\le\alpha_2$ and $\beta_1\le\beta_2$ the hierarchical version wins at every size. The
latency term falls from $2(gN-1)$ to $2(g-1)+2(N-1)$ steps, and the network bandwidth term falls by
about $g\times$ because only $1/g$ of the data crosses each NIC. With the figure's parameters
($g=N=8$) at 1\,GB: $T_\text{flat}\approx\SI{44.8}{\milli\second}$ and
$T_\text{hier}\approx\SI{9.3}{\milli\second}$, matching \code{llmsys.dist.costs}. Real NCCL achieves
much of this without an explicit hierarchy by running one ring per NIC.

\paragraph{Exercise 15.10}
70B: FSDP over 32 GPUs keeps $16\times70\times10^9/32=35$\,GB of persistent state per GPU, which fits in
80\,GB with room for transient units and activations. HSDP sharded over the 8 GPUs of a node needs
$16\times70\times10^9/8=140$\,GB per GPU and does not fit, so 70B must shard across nodes (or add TP/PP).
13B: HSDP needs $16\times13\times10^9/8=26$\,GB per GPU and fits. Its only cross-node traffic is an
all-reduce of each GPU's gradient shard (13B$/8$ parameters, about 3.3\,GB in bf16) across 4 replicas,
$2\cdot\frac34\cdot3.3\approx4.9$\,GB per GPU per step, overlapped with backward, as one collective
per bucket. FSDP over 32 GPUs would instead put two all-gathers and one reduce-scatter per
\emph{layer} on the network: more bytes and far more latency-sensitive messages. Use HSDP whenever
the model fits in one node's aggregate memory.

\paragraph{Exercise 15.11}
Expected shape of the results: per-parameter DDP is slowest when tensors are small and many, because
of latency; flat DDP fixes that but exposes all communication after backward; bucketed DDP hides
most of it once tokens per rank exceed $C/W$. Element-wise ZeRO-1 matches bucketed DDP's time to
within the extra copy and saves $12\Psi(1-1/p)$ bytes. FSDP without prefetch is slower by the sum of the
exposed all-gathers; with prefetch it approaches DDP once $T/n \gtrsim C/W$. Below that point FSDP's
extra $\Psi$ bytes per step show up as roughly $1.5\times$ DDP's communication time. The measured
crossover usually sits at somewhat more tokens than the model predicts, because achieved bus
bandwidth is below peak and kernels are less efficient at small batch.
